\chapter{Thesis Completion Roadmap}

ORCA is under advanced stages of development.
We first briefly describe the current state of the system, and then outline the remaining work
required to complete the thesis.

\section{Current State}

\paragraph{Completed Developmental Work}
Most core components of ORCA have been implemented.
The system supports overlay bootstrapping with multiple aggregators and topologies, telemetry
collection and forwarding mechanisms, and per-timestep analytics.
Observability requirements are programmed using a domain-specific language (DSL) called OrcaFlow:
the DSL enables defining transformations as SQL, applied in-situ while data is routed to
configurable \emph{sinks}.
The TS2PC protocol for the control path has been implemented.
End-to-end telemetry visualization is available via a DuckDB and a FlightSQL server embedded into
the controller, streamed to Grafana.
ORCA reuses the same infrastructure to export internal metrics for realtime visualization and
performance analysis.
Observability requirements from recent machine learning infrastructure have been reviewed and
incorporated in system semantics.

\paragraph{Validated Capabilities}
ORCA has been validated to be stable in a stress test with 128 MPI ranks loading a single
aggregator with 2--3 GB/s of telemetry.
Analytics were performed locally on the aggregator, achieving near-line rate performance under
certain conditions on 40 Gbps NICs.
Stress testing surfaced stability issues, which were mitigated by tuning pipeline parameters and
implementing flow control mechanisms.
We found that Arrow-based analytics had no trouble keeping up with the ingestion rate.
The RDMA-based transport layer, which is typically the most fragile part of such systems, is now
stable.

\paragraph{Remaining Developmental Work}
The primary remaining component is distributed query planning to fragment the OrcaFlow pipeline
across the overlay at operator-level granularity.
An initial prototype of the query planner exists but requires further refinement and integration.
Python bindings for Pytorch integration are also under development for the planned experiments.

\section{Evaluation Roadmap}

The evaluation will demonstrate the practical benefits of ORCA's relational approach in terms of
query performance, system scalability, and end-to-end impact.
The plan consists of the following experiments:

\begin{itemize}
  \item \textbf{Microbenchmark:}
        Compare query efficiency and analytical flexibility using on-disk Parquet telemetry versus
        traditional trace formats.
        AMR-inspired queries will serve as the HPC baseline, and Pytorch HTA will be used as the ML
        baseline.
  \item \textbf{Synthetic stress test:}
        Measure ORCA throughput and scalability as telemetry volume increases.
        Quantify application interference as a function of data load.
  \item \textbf{End-to-end workload study:}
        Measure runtime overhead and impact on scientific progress for a 4,096-rank AMR simulation and a
        data-parallel ML training job.
        Use tracing overheads as the baseline.
        Evaluate ORCA in two configurations: a lightweight in-situ aggregation setup and a full
        configuration that persists all telemetry at tracing volume.
\end{itemize}

Additionally, a survey of the machine learning observability literature is ongoing to assess
compatibility of representative analyses with ORCA's semantics.
This survey is not expected to have any experimental component.


% \section{Evaluation Roadmap}

% The evluation plan would involve demonstrating the analytical advantages of relational data
% formats over tracing formats, ORCA throughput and scalability and impact on application
% performance using synthetic stress tests, and impact on real HPC and ML workloads. Four
% experiments are planned:

% \begin{itemize} \item A microbenchmark to compare analytical efficiency of on-disk Parquet
% telemetry over tracing formats. Pytorch HTA (Holistic Trace Analysis) would be used as an ML
% baseline, and AMR-inspired queries as the HPC baseline. \item A synthetic stress test to evaluate
% ORCA throughput and scalability, and quantify interference as a function of data volume. \item A
% study of runtime overhead on two real workloads: a 4,096-rank AMR simulation and a data-parallel
% ML training job. The runtime overhead of tracing would be used as the baseline. ORCA would be run
% in two configurations: a lite config that does in-situ aggregation to compute summaries, and a
% full config that persists the same volume of data as full tracing. \end{itemize}

\section{Timeline}

The thesis is expected to be completed by the end of 2025.
The following timeline is proposed:

\begin{itemize}
  \item July--August 2025: Complete distributed query planning, Pytorch integration, and evaluation.
  \item September 2025: Paper and thesis writing, and job search.
  \item October 2025: Defend thesis.
\end{itemize}